\song{Marsyas a Apollón (Marsyas)}
% mara 19_11_06 Edit: Rádio 2022

\ve{}
\ch{A}Ta krásná dívka, \ch{Hmi}co se bojí o \ch{E}svoji krásu\\
\ch{A}Athéna \ch{F\guitarSharp{}mi}jméno má \ch{D}za starých \ch{A}dávných časů\\
\ch{A}Odhodí flétnu, \ch{Hmi}hrát nejde s \ch{E}nehybnou tváří\\
\ch{A}Ten, kdo ji \ch{F\guitarSharp{}mi}najde dřív, \ch{D}tomu se \ch{A}přání zmaří\ch{A/As}

\cho{}
\ch{F\guitarSharp{}mi}Tak si Marsyas \ch{E}mámen flétnou \ch{F\guitarSharp{}mi}věří, že \ch{D}musí přetnout\\
/: \ch{A}jedno pravidlo, \ch{E}sázku a \ch{D}hrát \ch{A}líp než \ch{E}bůh :/

\ve{}
Bláznivý nápad, snad nejvýš Marsyas míří\\
Apollón souhlasí, oba se s trestem smíří\\
Král Midas má říct, kdo je lepší, Apollón zpívá\\
O život soupeří, jen jeden vítěz bývá

\cho{}

\ve{}
Obrátí nástroj, už ví, že nebude chválen\\
Prohrál a zápolí podveden vůlí krále\\
Sám v tichém hloučku, sám na strom připraví ráhno\\
Satyra k hrůze všech zaživa z kůže stáhnou

\cho{}
Tak si Apollón změřil síly, každý se musel mýlit\\
/: Nikdo nemůže kouzlit a hrát líp než bůh :/

\ve{}
Ta krásná dívka, co se bála o svoji krásu\\
Dárkyně moudrosti za starých dávných časů\\
Teď v tichém hloučku, v jejích rukou úroda, spása\\
Athéna jméno má, chybí jí tvář a krása

\cho{}
Dy dy dy\ldots{}

\esong{}